\documentclass[11pt]{article}

\usepackage[margin=1in]{geometry}
\usepackage{booktabs}
\usepackage{hyperref}
\usepackage{amsmath}

\title{Luxembourg and Timor-Leste Case Studies Recreated With the Scalable Distance Pipeline}
\author{Public Infrastructure Service Access}
\date{}

\begin{document}
\maketitle

\section{Purpose}

This note restates the Luxembourg school-access and Timor-Leste health-facility
case studies using the reengineered
\texttt{scalable\_general\_distances\_per\_country} pipeline.  The goal is to
describe the same analytical narrative as the original paper, but in terms of
the new modular contracts: shared cache manifests, country source descriptors,
source and target tables, interchangeable routing strategies, sparse distance
matrices, and optimization-ready coverage neighborhoods.

The companion notebook
\texttt{luxembourg\_api\_illustration.ipynb} contains executable miniature
versions of the cases.  Those examples are intentionally small enough to run on
a fresh development machine, while the final notebook cells show the country
scale API and CLI calls that use the shared cache.

\section{Pipeline Architecture Used by Both Cases}

Both case studies use the same sequence of self-contained modules.

\begin{enumerate}
  \item \textbf{Country source resolution.}
  \texttt{CountryDataSources} resolves the Geofabrik PBF filename, WorldPop
  filename, download URLs, and deterministic cache paths before any data are
  downloaded.

  \item \textbf{Shared cache and manifest.}
  Downloaded artefacts are stored under the shared cache used by the original
  implementation.  The cache layout and manifest make it explicit which files
  are downloaded, generated, or reused.

  \item \textbf{Data extraction.}
  The production runner extracts a drivable network from the OSM PBF, converts
  WorldPop raster cells to target points, and extracts OSM facilities matching
  the requested amenity values.

  \item \textbf{Routing strategy.}
  The router is selected by name.  \texttt{NetworkXRouter} is the portable
  default.  \texttt{PandanaRouter} is imported only when selected, so deployment
  environments without a compatible Pandana and NumPy stack can still use the
  package.

  \item \textbf{Distance matrix output.}
  The same source and target tables feed the router, and the matrix writer can
  produce combined, split, or both output layouts.  The split layout supports
  large-country runs by keeping source-target distance chunks lightweight.

  \item \textbf{Downstream analysis.}
  Descriptive coverage and prescriptive maximum-covering models consume the
  sparse matrix without depending on the implementation details of the routing
  backend.
\end{enumerate}

\section{Case Definitions}

\begin{table}[h]
\centering
\begin{tabular}{p{0.22\linewidth}p{0.34\linewidth}p{0.34\linewidth}}
\toprule
Case & Paper settings & Paper-scale reported result \\
\midrule
Luxembourg schools &
OSM \texttt{amenity=school}; no population aggregation; retained matrix
threshold of 1,000 m; context map saved. &
394 OSM school features; 102,061 population targets representing 478,010
people; 266,225 people within 1,000 m of an OSM school. \\
\addlinespace
Timor-Leste health &
Default health amenities plus generated 5 km candidate sites; aggregate factor
8; retained matrix threshold of 10 km; maximum covering with \(p=20\). &
28,162 population targets; 860 sources; 100,997 retained distance rows; Gurobi
and HiGHS both cover 757,053 represented people at 10 km. \\
\bottomrule
\end{tabular}
\caption{Paper case definitions restated as scalable-pipeline inputs.}
\end{table}

\section{Luxembourg School Access}

The Luxembourg case is descriptive.  OSM schools are treated as existing source
locations and WorldPop cells are treated as population targets.  The reengineered
pipeline resolves the country data as:

\begin{verbatim}
CountryDataSources(
    country_slug="luxembourg",
    iso3="LUX",
    base_dir=shared_cache / "luxembourg_data",
    worldpop_filename="lux_ppp_2020.tif",
)
\end{verbatim}

The production configuration is:

\begin{verbatim}
ProductionRunConfig(
    sources=luxembourg_sources,
    output_dir=shared_cache / "luxembourg_data" / "outputs",
    run_tag="luxembourg_schools_networkx",
    amenity_values=("school",),
    router="networkx",
    matrix_output_mode="both",
)
\end{verbatim}

For each target \(i\), the analysis asks whether at least one school \(j\) is
within a road-network threshold \(R\):

\[
  \operatorname{covered}_i(R) =
  \begin{cases}
    1 & \text{if } \min_{j \in J_{\mathrm{school}}} d_{ij} \le R,\\
    0 & \text{otherwise.}
  \end{cases}
\]

The population covered at threshold \(R\) is then
\[
  \sum_i w_i \operatorname{covered}_i(R),
\]
where \(w_i\) is the represented population of target \(i\).  A target is counted
only once per threshold even if several schools are reachable.  For school
ranking, the same matrix can assign each covered target to its nearest school.

The corresponding CLI call is:

\begin{verbatim}
py -m scalable_distances.cli run ^
  --country-slug luxembourg ^
  --iso3 LUX ^
  --base-dir C:\local\Download_Depot\luxembourg_data ^
  --output-dir C:\local\Download_Depot\luxembourg_data\outputs ^
  --run-tag luxembourg_schools_networkx ^
  --amenity school ^
  --router networkx ^
  --matrix-output-mode both
\end{verbatim}

\section{Timor-Leste Health Facility Coverage}

The Timor-Leste case is prescriptive.  It builds a sparse matrix between
population targets and health-related source locations, then chooses a limited
number of sources to maximize covered population.  The reengineered pipeline
resolves the country data as:

\begin{verbatim}
CountryDataSources(
    country_slug="east-timor",
    iso3="TLS",
    base_dir=shared_cache / "east-timor_data",
    worldpop_filename="tls_ppp_2020.tif",
)
\end{verbatim}

The health amenity run is configured as:

\begin{verbatim}
ProductionRunConfig(
    sources=timor_leste_sources,
    output_dir=shared_cache / "east-timor_data" / "outputs",
    run_tag="timor_leste_health_networkx",
    amenity_values=(
        "hospital", "clinic", "doctors", "dentist", "pharmacy",
        "health_post", "nursing_home", "social_facility",
    ),
    router="networkx",
    aggregate_factor=8,
    matrix_output_mode="both",
)
\end{verbatim}

The paper case also includes generated 5 km candidate sites.  In the new design,
candidate generation belongs to the same source-table contract as OSM amenities:
candidate points can be added as an additional source layer before snapping and
routing, while the router and matrix writer remain unchanged.

Let \(N(i)\) be the set of source locations within 10 km of target \(i\), based
on the sparse routing matrix.  The maximum-covering model is:

\[
\max \sum_i w_i z_i
\]
\[
z_i \le \sum_{j \in N(i)} y_j \quad \forall i,
\qquad
\sum_j y_j \le p,
\qquad
z_i, y_j \in \{0, 1\}.
\]

Here \(y_j=1\) if source \(j\) is selected and \(z_i=1\) if population target
\(i\) is covered.  In the paper-scale case, \(p=20\), the coverage radius is
10 km, and both Gurobi and HiGHS reach the same covered population.

The corresponding CLI scaffold is:

\begin{verbatim}
py -m scalable_distances.cli run ^
  --country-slug east-timor ^
  --iso3 TLS ^
  --base-dir C:\local\Download_Depot\east-timor_data ^
  --output-dir C:\local\Download_Depot\east-timor_data\outputs ^
  --run-tag timor_leste_health_networkx ^
  --amenity hospital --amenity clinic --amenity doctors ^
  --amenity dentist --amenity pharmacy --amenity health_post ^
  --amenity nursing_home --amenity social_facility ^
  --router networkx ^
  --aggregate-factor 8 ^
  --matrix-output-mode both
\end{verbatim}

\section{Reproducibility Checks}

The new pipeline makes the following checks explicit for both cases:

\begin{itemize}
  \item source-data filenames and URLs are resolved before download;
  \item downloads are cached in the shared cache and reused on repeated runs;
  \item source and target tables expose stable IDs, geometry, and attributes;
  \item route computation records router name, network size, source count,
  target count, and matrix row count;
  \item split and combined matrix outputs can be compared from the same run;
  \item Pandana is not imported unless the Pandana strategy is requested.
\end{itemize}

These checks allow the original implementation and the scalable implementation
to be compared at the level of cache contents, road-network inputs, coordinate
handling, snapping, origin-destination pairs, distance calculations, and final
summary statistics.

\section{Interpretation}

The core scientific narrative is unchanged: the Luxembourg case measures the
population already close to schools, while the Timor-Leste case uses the same
distance-matrix representation to support facility-location decisions.  The
engineering narrative changes substantially.  Instead of one monolithic script,
the reengineered version exposes lightweight data contracts, lazy routing
strategies, explicit persistence choices, and cache metadata that can be reused
by notebooks, command-line runs, tests, and future geocoding or optimization
modules.

\end{document}
